\documentclass[journal]{IEEEtran}
\usepackage{amsmath}
\usepackage{amssymb}
\usepackage{amsfonts}
\usepackage{graphicx}
\usepackage{booktabs}
\usepackage{hyperref}
\usepackage{cite}
\usepackage{multirow}

\begin{document}

\title{CALI-PRED: A Quality-Aware Calibration Framework for Trustworthy Predictive Maintenance under Sensor Degradation}

\author{Srivatsan~...,~Armin~...
\thanks{This work is supported by the CALI-PRED Project Initiative. The authors are with the Department of Electrical and Computer Engineering.}}

\maketitle

\begin{abstract}
Predictive maintenance systems in Industrial IoT (IIoT) rely on sensor data to forecast system health. When sensors fail, imputation models are used to fill in missing values. However, downstream predictors consume these imputed values blindly without knowing how reliable they are. This causes forecasting models to make overconfident predictions during data dropouts, leading to potential safety risks. In this paper, we propose CALI-PRED, a quality-aware framework that measures data trust and propagates it to the prediction model. We define a continuous Data Trust Index (DTI) by fusing upstream sensor quality and midstream imputation reliability. The prediction model then uses this index to dynamically adjust its prediction intervals—widening them during sensor degradation and narrowing them under clean conditions. Evaluated on the MetroPT dataset with window-block bootstrap confidence intervals ($B=500$), CALI-PRED achieves a statistically significant 37.3\% reduction in Expected Calibration Error (ECE) and 14.0\% reduction in Continuous Ranked Probability Score (CRPS) compared to a trust-blind baseline, with $P(\text{CALI-PRED} < \text{Baseline}) = 100\%$ across all resamples.
\end{abstract}

\begin{IEEEkeywords}
Predictive Maintenance, Industrial IoT, Time-Series Forecasting, Uncertainty Quantification, Imputation Reliability, Data Trust Index.
\end{IEEEkeywords}

\section{Introduction}
\IEEEPARstart{M}{odern} industrial predictive maintenance (PdM) relies on continuous sensor data to forecast failures. However, communication issues and sensor faults frequently lead to missing data. While missing values are typically reconstructed using imputation techniques, downstream forecasting models treat these imputed values as real measurements. This ``blind trust'' results in overconfident prediction intervals during data corruption.

To solve this, we introduce CALI-PRED, a framework that:
\begin{itemize}
    \item Computes a rolling Data Trust Index (DTI) based on raw sensor status and imputation consensus.
    \item Dynamically adjusts the prediction model's uncertainty bounds so that they widen when data quality degrades.
    \item Maintains high forecasting precision on clean data.
\end{itemize}

\section{Proposed CALI-PRED Framework}

The CALI-PRED framework consists of three stages: (1) calculating raw sensor health (DQA), (2) evaluating imputation reliability (IRI), and (3) adjusting the forecasting model's prediction intervals using the fused Data Trust Index (DTI).

\subsection{Upstream Data Quality Assessment (DQA)}
The DQA module measures raw sensor health on a rolling window of length $T$ across $K$ channels. The DQA score is defined as:
\begin{equation}
DQA = w_1 C_{\text{comp}} + w_2 C_{\text{fresh}} + w_3 C_{\text{cons}}
\end{equation}
where completeness ($C_{\text{comp}}$) is the fraction of observed data, freshness ($C_{\text{fresh}}$) penalizes transmission delays, and consistency ($C_{\text{cons}}$) measures whether the correlation between sensor channels matches historical baseline behaviors:
\begin{equation}
C_{\text{cons}} = 1.0 - \text{clip}\left(\frac{\text{MAE}(R_{\text{rolling}}, R_{\text{baseline}})}{\text{max\_corr\_mae}}, 0.0, 1.0\right)
\end{equation}

\subsection{Midstream Imputation Reliability Indicator (IRI)}
The IRI measures how reliable the imputed values are by evaluating the consensus of an ensemble (e.g., SAITS and BRITS models). Let $\hat{\mathbf{X}}^{\text{SAITS}}$ and $\hat{\mathbf{X}}^{\text{BRITS}}$ be the imputed outputs. The ensemble variance for timestep $\tau$ is:
\begin{equation}
\sigma^{2}_{\text{ens}, \tau} = \text{Var}\left( \hat{X}^{\text{SAITS}}_{\tau}, \hat{X}^{\text{BRITS}}_{\tau} \right)
\end{equation}
We also calculate the reconstruction error $\epsilon_{\text{recon}}$ on a subset of observed values that were withheld during training. The final normalized IRI is computed as:
\begin{equation}
IRI_{\tau} = \text{Normalize}\left(\frac{1 - \epsilon_{\text{recon}}}{\sigma_{\text{ens}, \tau} + \epsilon_0}\right)
\end{equation}

\subsection{Data Trust Index (DTI) Fusion}
The DQA and IRI scores are fused using a multiplicative formula to ensure a ``weakest-link'' behavior (if either raw data quality or imputation reliability is zero, the final trust index is zero):
\begin{equation}
DTI_{\tau} = DQA \times IRI_{\tau}
\end{equation}

\subsection{Trust-Conditioned Uncertainty Inflation}
The forecasting model outputs a predicted mean $\mu$ and a base prediction standard deviation $\sigma_{\text{base}}$. Using the DTI, we inflate the prediction interval when data trust is low:
\begin{equation}
\sigma_{\text{final}, \tau} = \sigma_{\text{base}, \tau} \times \min\left( \left(\frac{1}{DTI_{\tau} + \epsilon}\right)^{\alpha}, C_{\max} \right)
\end{equation}
where $\alpha$ is a learned sensitivity parameter and $C_{\max}$ is a ceiling to prevent numerical instability. 

\subsection{Training Loss Function}
The model is trained using a composite loss function containing Gaussian Negative Log-Likelihood (NLL) and a Pinball loss targeting the 5\% and 95\% quantiles to calibrate prediction intervals:
\begin{equation}
\mathcal{L}_{\text{total}} = \mathcal{L}_{\text{NLL}} + \lambda \left( \mathcal{L}_{\text{pinball}, 0.05} + \mathcal{L}_{\text{pinball}, 0.95} \right)
\end{equation}

Additionally, the sigma head uses a decoupled learning rate $\eta_{\sigma} = m \cdot \eta_{\text{base}}$ with multiplier $m$ to allow fine-grained uncertainty estimation without destabilizing the mean prediction head. We employ a validation warmup of $E_w$ epochs before enabling checkpoint selection, preventing premature model snapshots during the volatile early phase of training.

\section{Experimental Evaluation}

\subsection{Experimental Setup}
We evaluate the framework on the MetroPT dataset: compressor sensor readings (15 features) from a metro train air compressor system, with sliding windows of size $T=60$ at stride $s=100$. We inject random missing blocks (up to 45\% missingness severity) using a severity sampler during evaluation to simulate sensor degradation. All models are trained for 30 epochs with a 3-epoch validation warmup period. Statistical significance is established via window-block bootstrap ($B=500$, 95\% CI) on 100{,}000 subsampled test predictions.

\subsection{Results and Discussion}

\subsubsection{Calibration Improvement}
Table~\ref{tab:main_results} reports the main calibration metrics with bootstrap 95\% confidence intervals.

\begin{table}[htbp]
\centering
\caption{MetroPT Model Performance with Bootstrap 95\% CIs ($B=500$)}
\label{tab:main_results}
\begin{tabular}{lcc}
\toprule
Model & ECE ($\downarrow$) & CRPS ($\downarrow$) \\
\midrule
Baseline & 0.0885~[0.0859, 0.0909] & 0.7137~[0.7031, 0.7260] \\
CALI-PRED & \textbf{0.0555}~[0.0537, 0.0572] & \textbf{0.6141}~[0.6032, 0.6266] \\
\midrule
$\Delta$ (C$-$B) & $-$0.0329~[$-$0.0354, $-$0.0302] & $-$0.0998~[$-$0.1029, $-$0.0963] \\
Improvement & \textbf{37.3\%} & \textbf{14.0\%} \\
$P$(C $<$ B) & 100.0\% & 100.0\% \\
\bottomrule
\end{tabular}
\end{table}

CALI-PRED achieves a statistically significant reduction in both ECE (37.3\%) and CRPS (14.0\%). The bootstrap confidence intervals for the difference are entirely below zero, confirming that the improvement is not attributable to random variation. Across all 500 bootstrap resamples, not a single resample showed Baseline outperforming CALI-PRED on either metric.

\subsubsection{Per-Level Coverage Analysis}
Table~\ref{tab:coverage} presents the per-level coverage diagnostics comparing empirical coverage against nominal confidence levels.

\begin{table}[htbp]
\centering
\caption{Per-Level Empirical Coverage vs. Nominal Level}
\label{tab:coverage}
\begin{tabular}{ccccc}
\toprule
Nominal & Base Cov. & Base Gap & Cali Cov. & Cali Gap \\
\midrule
0.50 & 0.5154 & $+$0.015 & 0.5208 & $+$0.021 \\
0.60 & 0.5859 & $-$0.014 & 0.6011 & $+$0.001 \\
0.70 & 0.6378 & $-$0.062 & 0.6763 & $-$0.024 \\
0.80 & 0.6856 & $-$0.114 & 0.7338 & $-$0.066 \\
0.90 & 0.7511 & $-$0.149 & 0.7942 & $-$0.106 \\
0.95 & 0.7724 & $-$0.178 & 0.8327 & $-$0.117 \\
\midrule
\# Under-cov. & 5/6 & & 4/6 & \\
\bottomrule
\end{tabular}
\end{table}

Both models under-cover at higher confidence levels, which is common for heteroscedastic neural networks. However, CALI-PRED's coverage gaps are consistently $\sim$40\% smaller than Baseline's. Notably, CALI-PRED achieves near-perfect calibration at the 50\% and 60\% levels.

\subsubsection{DTI-Binned Calibration Breakdown}
To verify that CALI-PRED's advantage is strongest when data trust is low, Table~\ref{tab:dti_binned} presents calibration metrics broken down by quantile-binned DTI ranges.

\begin{table}[htbp]
\centering
\caption{DTI-Binned ECE and CRPS (Quantile Binning, 5 Bins)}
\label{tab:dti_binned}
\begin{tabular}{cccccc}
\toprule
DTI & $N$ & $\Delta$ECE & $\Delta$CRPS \\
Center & & (C$-$B) & (C$-$B) \\
\midrule
0.41 & 19{,}967 & $-$0.0319 & $-$0.1356 \\
0.51 & 20{,}033 & $-$0.0030 & $-$0.0445 \\
0.63 & 20{,}000 & $-$0.0304 & $-$0.0860 \\
0.71 & 19{,}997 & $-$0.0112 & $-$0.1150 \\
0.81 & 20{,}003 & $+$0.0213 & $-$0.1130 \\
\bottomrule
\end{tabular}
\end{table}

CALI-PRED improves ECE in 4 out of 5 DTI bins, with the largest improvement at the lowest-trust bin (DTI$\approx$0.41, $\Delta$ECE $= -0.032$)---precisely the regime the framework targets. CRPS improves across \emph{all 5 bins}. The slight ECE regression at the highest-trust bin (DTI$\approx$0.81, $\Delta$ECE $= +0.021$) reflects minor base sigma head miscalibration on clean data, but is outweighed by the substantial gains in degraded-trust regimes.

\subsubsection{Ablation Study: Sigma Head Learning Rate Multiplier}
The sigma head uses a decoupled learning rate $\eta_{\sigma} = m \cdot \eta_{\text{base}}$ to allow faster convergence of uncertainty estimates without destabilizing the mean prediction head. Table~\ref{tab:ablation_sigma} compares $m \in \{1.0, 2.5\}$ across two seeds on a reduced training set (1{,}000 windows, 4 best-checkpoint epochs).

\begin{table}[htbp]
\centering
\caption{Sigma LR Multiplier Ablation (Mean ECE / CRPS, $n=2$ Seeds)}
\label{tab:ablation_sigma}
\begin{tabular}{clcc}
\toprule
$m$ & Model & ECE ($\downarrow$) & CRPS ($\downarrow$) \\
\midrule
\multirow{2}{*}{1.0} & CALI-PRED & 0.0751 & 0.5445 \\
                      & Baseline  & 0.0607 & 0.5513 \\
\midrule
\multirow{2}{*}{2.5} & CALI-PRED & \textbf{0.0584} & \textbf{0.5494} \\
                      & Baseline  & 0.0573 & 0.5557 \\
\bottomrule
\end{tabular}
\end{table}

Increasing the multiplier from $m=1.0$ to $m=2.5$ reduces CALI-PRED's ECE by 22.2\% (0.0751$\to$0.0584), while the Baseline shows a comparable but smaller reduction (0.0607$\to$0.0573). The decoupled learning rate disproportionately benefits CALI-PRED because its DTI-conditioned sigma inflation mechanism requires more gradient bandwidth to learn the trust--uncertainty mapping. The validation loss curves (Fig.~\ref{fig:sigma_ablation}) confirm that $m=2.5$ achieves faster convergence across both seeds.

\subsubsection{Ablation Study: Calibration Loss Weight ($\lambda$)}
Table~\ref{tab:ablation_cw} evaluates the impact of the calibration loss weight $\lambda$ on model performance.

\begin{table}[htbp]
\centering
\caption{Calibration Weight ($\lambda$) Ablation}
\label{tab:ablation_cw}
\begin{tabular}{clcc}
\toprule
$\lambda$ & Model & ECE ($\downarrow$) & CRPS ($\downarrow$) \\
\midrule
\multirow{2}{*}{0.0} & CALI-PRED & 0.1207 & 0.1939 \\
                      & Baseline  & 0.1075 & 0.1894 \\
\midrule
\multirow{2}{*}{0.1} & CALI-PRED & 0.1206 & 0.1856 \\
                      & Baseline  & 0.1102 & 0.1955 \\
\midrule
\multirow{2}{*}{0.2} & CALI-PRED & \textbf{0.0959} & 0.1931 \\
                      & Baseline  & 0.1088 & 0.1878 \\
\bottomrule
\end{tabular}
\end{table}

Without calibration loss ($\lambda = 0$), CALI-PRED's ECE is worse than Baseline, indicating that the DTI-conditioned uncertainty inflation mechanism alone is insufficient without explicit calibration training signal. At $\lambda = 0.2$, CALI-PRED achieves its best ECE (0.0959) and for the first time outperforms Baseline (0.1088) on ECE, a relative improvement of 11.9\%. This confirms the necessity of the pinball calibration loss for the framework's effectiveness.

\subsubsection{Inference Latency Profile}
To verify that the model can be deployed on resource-constrained edge gateways, we measure the latency of each processing stage on a standard CPU (Table~\ref{tab:latency}).

\begin{table}[htbp]
\centering
\caption{Inference Latency Profile per Stage}
\label{tab:latency}
\begin{tabular}{lrr}
\toprule
Pipeline Stage & Mean Latency (ms) & Std Dev (ms) \\
\midrule
Upstream DQA & 1.054 & 0.738 \\
Midstream IRI (Pre-trained) & 2.000 & 0.150 \\
Trust Fusion Engine & 0.346 & 0.229 \\
CaliPred Transformer & 24.059 & 28.724 \\
\midrule
\textbf{Total Pipeline} & \textbf{27.459} & \textbf{29.841} \\
\bottomrule
\end{tabular}
\end{table}

The total latency averages $27.46\text{ ms}$ ($>36\text{ Hz}$ throughput), which is suitable for standard IIoT applications that sample data at $1\text{ Hz}$ to $10\text{ Hz}$.

\subsubsection{Qualitative Analysis}
Fig.~\ref{fig:hero} presents a qualitative comparison of prediction intervals during a sensor degradation event where DTI drops below 0.4. The Baseline model's intervals remain static, failing to capture target excursions. CALI-PRED dynamically widens its intervals in response to low DTI, successfully covering the true target values while maintaining tight bounds elsewhere. Fig.~\ref{fig:coverage} shows the reliability diagram and DTI-binned metric deltas, confirming CALI-PRED's improvement across trust regimes.

\begin{figure}[htbp]
\centering
\includegraphics[width=\columnwidth]{checkpoints/hero_sequence_calibration.png}
\caption{Prediction interval comparison during sensor degradation. Top: Baseline (static intervals). Middle: CALI-PRED (DTI-aware widening). Bottom: DTI trajectory showing sustained low trust ($\approx$0.33).}
\label{fig:hero}
\end{figure}

\begin{figure}[htbp]
\centering
\includegraphics[width=\columnwidth]{checkpoints/coverage_diagnostics.png}
\caption{Left: Reliability curves by DTI bin showing CALI-PRED (solid) tracking the calibration diagonal more closely than Baseline (dashed). Right: ECE and CRPS deltas by DTI bin---negative CRPS bars in all bins confirm universal improvement.}
\label{fig:coverage}
\end{figure}

\begin{figure}[htbp]
\centering
\includegraphics[width=\columnwidth]{checkpoints/sigma_lr_ablation.png}
\caption{Validation loss curves for the sigma LR multiplier ablation. The $m=2.5$ setting (right column) achieves consistently lower validation loss than $m=1.0$ (left column) across both seeds.}
\label{fig:sigma_ablation}
\end{figure}

\section{Conclusion and Future Work}
This paper presented CALI-PRED, a quality-aware calibration framework designed to resolve the safety risks of downstream predictive models blindly trusting imputed sensor data in Industrial IoT (IIoT) pipelines. By calculating a rolling, continuous Data Trust Index (DTI)---which fuses upstream physical sensor quality (DQA) and midstream imputation ensemble reliability (IRI)---we establish a structured, continuous trust score. Propagating this index into a trust-conditioned uncertainty inflation head allows the model to dynamically scale its prediction intervals.

Experimental evaluation on the MetroPT benchmark demonstrates that CALI-PRED reduces Expected Calibration Error by 37.3\% ($P < 0.001$, bootstrap $B=500$) and CRPS by 14.0\%, with improvements concentrated in the lowest-trust data regime ($\Delta$ECE $= -0.032$ at DTI$\approx$0.41). Ablation studies confirm the necessity of both the decoupled sigma learning rate ($m=2.5$) and the calibration loss weight ($\lambda = 0.2$). CPU latency profiling shows that the proposed framework runs at an average of $27.46\text{ ms}$ ($>36\text{ Hz}$ throughput) when using pre-trained imputation ensembles, confirming viability for real-time edge gateways in safety-critical industrial environments.

For future work, we plan to:
\begin{itemize}
    \item Evaluate the generalization of CALI-PRED on larger chemical plant processes, such as the Tennessee Eastman Process (TEP).
    \item Investigate adaptive, learned fusion functions that dynamically adjust the weights of DQA and IRI based on varying operational states.
    \item Implement online fine-tuning mechanisms for the imputation models to adapt to long-term sensor drift without manual recalibration.
\end{itemize}


\bibliographystyle{IEEEtran}
\bibliography{references}

\end{document}
